\begin{table}[!htbp]
\centering
\caption{A fixed cohort decomposition through twenty days after implementation}\label{tab:horizon}
\small
\setlength{\tabcolsep}{5pt}
\begin{tabular}{llrrrrlr}
\toprule
Type & Period & $N$ & Before ED & After ED & Cumulative & CR1 interval & $p_{WCR}$ \\
\midrule
Addition & 2017--2026 & 169 & 2.43 & 0.87 & 3.30 & [1.31, 5.28] & $<0.001$ \\
Addition & 2022--2026 & 87 & 1.77 & -0.06 & 1.71 & [0.72, 2.69] & 0.010 \\
Deletion & 2017--2026 & 136 & -0.70 & -0.38 & -1.07 & [-3.68, 1.54] & 0.414 \\
Deletion & 2022--2026 & 75 & -0.93 & -0.95 & -1.88 & [-5.72, 1.96] & 0.373 \\
\bottomrule
\end{tabular}
\begin{minipage}{\linewidth}\footnotesize\vspace{5pt}
Exploratory extension. Before ED is CAR $[AD-1,ED-1]$; After ED is CAR $[ED,ED+20]$; Cumulative is their event-level sum, CAR $[AD-1,ED+20]$. Every entry in a row uses exactly the same complete-case events. All values are percentages. The interval is analytical CR1 with $t_{G-1}$; $p_{WCR}$ is a separately computed restricted wild cluster bootstrap $p$-value using 49,999 draws and the finite-simulation correction $(k+1)/(B+1)$. These are distinct inferential summaries. No claim of permanent or causal price impact follows from a nonzero finite-horizon cumulative return. Source: Authors' calculations; exact inputs are listed in the replication package.
\end{minipage}
\end{table}
